\Article{231}
Die Gesamtheit \textit{aller} Einzelcharaktere einer gegebenen Form oder Klasse bildet den Gesamtcharakter dieser Form oder Klasse. So ist \textit{zum Beispiel} der Gesamtcharakter der Form $(10,3,17)$ oder der ganzen von ihr repräsentierten Klasse $1,4$; $N7$; $N23$. Ebenso ist der Gesamtcharakter der Form $(7,1,-17)$ $7,8$; $R3$; $N5$, denn der Einzelcharakter $3,4$ wird hier weggelassen, weil er bereits im Charakter $7,8$ enthalten ist. --- Darauf gründen wir die Einteilung der ganzen Ordnung der eigentlich primitiven Klassen (der positiven, wenn die Determinante negativ ist) mit gegebener Determinante in mehrere verschiedene \textit{Geschlechter}: Alle Klassen mit demselben Gesamtcharakter gehören zum selben Geschlecht, Klassen mit verschiedenen Gesamtcharakteren dagegen zu verschiedenen Geschlechtern. Jedem Geschlecht geben wir den Gesamtcharakter der darin enthaltenen Klassen. So gibt es \textit{zum Beispiel} für die Determinante $-161$ sechzehn
\SourceLeaf{239}
positive eigentlich primitive Klassen. Sie verteilen sich auf folgende vier Geschlechter:
\begin{center}
\begin{tabular}{c|l}
Charakter & \multicolumn{1}{c}{Repräsentierende Formen der Klassen} \\
$1,4;\ R7;\ R23$ & $(1,0,161)$, $(2,1,81)$, $(9,1,18)$, $(9,-1,18)$ \\
$1,4;\ N7;\ N23$ & $(5,2,33)$, $(5,-2,33)$, $(10,3,17)$, $(10,-3,17)$ \\
$3,4;\ R7;\ N23$ & $(7,0,23)$, $(11,2,15)$, $(11,-2,15)$, $(14,7,15)$ \\
$3,4;\ N7;\ R23$ & $(3,1,54)$, $(3,-1,54)$, $(6,1,27)$, $(6,-1,27)$
\end{tabular}
\end{center}

Für die Anzahl der verschiedenen Gesamtcharaktere, die von vornherein möglich sind, gilt Folgendes.

I. Ist die Determinante $D$ durch $8$ teilbar, so sind bezüglich $8$ vier verschiedene Einzelcharaktere möglich; die Zahl $4$ liefert keinen eigenen Charakter (Anmerkung zum vorigen Artikel). Außerdem gibt es bezüglich jedes ungeraden Primteilers von $D$ zwei Charaktere. Ist deren Anzahl $m$, so gibt es insgesamt $2^{m+2}$ verschiedene Gesamtcharaktere (wobei wir $m=0$ setzen, wenn $D$ eine Zweierpotenz ist).

II. Ist die Determinante $D$ nicht durch $8$ teilbar, wohl aber durch $4$ und außerdem durch $m$ ungerade Primzahlen, so gibt es insgesamt $2^{m+1}$ verschiedene Gesamtcharaktere.

III. Ist die Determinante $D$ gerade, aber nicht durch $4$ teilbar, so ist sie entweder $\equiv2\ (\mathrm{mod.}\ 8)$ oder $\equiv6$. Im ersten Fall gibt es bezüglich $8$ zwei Einzelcharaktere, nämlich $1\ \mathit{und}\ 7,8$ und $3\ \mathit{und}\ 5,8$; im zweiten ebenso viele. Setzen wir also die Anzahl der ungeraden Primteiler von $D$ $=m$, so gibt es insgesamt $2^{m+1}$ verschiedene Gesamtcharaktere.

IV. Ist $D$ ungerade, so ist sie entweder $\equiv1$ oder $\equiv3\ (\mathrm{mod.}\ 4)$. Im zweiten Fall gibt es bezüglich $4$ zwei verschiedene Charaktere; im ersten geht eine solche Beziehung nicht in den Gesamtcharakter ein. Bezeichnet $m$ also dasselbe wie zuvor, so gibt es im ersten Fall $2^m$, im zweiten $2^{m+1}$ verschiedene Gesamtcharaktere.

Es ist aber gut zu beachten, dass daraus keineswegs folgt, dass tatsächlich ebenso viele Geschlechter
\SourceLeaf{240}
existieren, wie verschiedene Charaktere von vornherein möglich sind. In unserem Beispiel entsprechen nur der Hälfte von ihnen tatsächlich Klassen oder Geschlechter. Es gibt keine positiven Klassen mit den Charakteren $1,4;\ R7;\ N23$ oder $1,4;\ N7;\ R23$; oder $3,4;\ R7;\ R23$ oder $3,4;\ N7;\ N23$. Auf diesen sehr wichtigen Gegenstand gehen wir unten ausführlicher ein.

Die Form $(1,0,-D)$, die zweifellos unter allen Formen mit der Determinante $D$ als die einfachste anzusehen ist, nennen wir von nun an \textit{Hauptform}. Die ganze Klasse, die sie enthält, heißt \textit{Hauptklasse}, und das ganze Geschlecht, das die Hauptklasse enthält, \textit{Hauptgeschlecht}. Wir müssen daher die Hauptform, eine Form aus der Hauptklasse und eine Form aus dem Hauptgeschlecht sorgfältig unterscheiden; ebenso die Hauptklasse und eine Klasse aus dem Hauptgeschlecht. Diese Bezeichnungen verwenden wir stets, auch wenn es für eine Determinante keine anderen Klassen als die Hauptklasse oder keine anderen Geschlechter als das Hauptgeschlecht gibt. Das geschieht \textit{zum Beispiel} häufig, wenn $D$ eine positive Primzahl der Form $4n+1$ ist.

\Needspace{8\baselineskip}
\Article{232}
Zwar haben wir die Charaktere der Formen zunächst erläutert, um darauf die Unterteilung der \textit{positiven eigentlich primitiven} Ordnung zu gründen. Doch können wir dasselbe auch auf negative Formen und Klassen oder auf uneigentlich primitive anwenden. Nach demselben Grundsatz lassen sich sowohl die positive uneigentlich primitive Ordnung als auch die negative eigentlich primitive und die negative uneigentlich primitive Ordnung in Geschlechter einteilen. So kann \textit{zum Beispiel}, nachdem die eigentlich primitive Ordnung der Formen mit der Determinante $145$ in die folgenden beiden Geschlechter eingeteilt ist,
\begin{center}
\begin{tabular}{l|l}
$R5,\ R29$ & $(1,0,-145)$, $(5,0,-29)$ \\
$N5,\ N29$ & $(3,1,-48)$, $(3,-1,-48)$
\end{tabular}
\end{center}
auch die uneigentlich primitive Ordnung ebenso in zwei Geschlechter eingeteilt werden:
\begin{center}
\begin{tabular}{l|l}
$R5,\ R29$ & $(4,1,-36)$, $(4,-1,-36)$ \\
$N5,\ N29$ & $(2,1,-72)$, $(10,5,-12)$
\end{tabular}
\end{center}
Oder, wie die positiven Klassen von Formen mit der Determinante $-129$ in vier Geschlechter zerfallen,
\SourceLeaf{241}
\begin{center}
\begin{tabular}{l|l}
$1,4;\ R3;\ R43$ & $(1,0,129)$, $(10,1,13)$, $(10,-1,13)$ \\
$1,4;\ N3;\ N43$ & $(2,1,65)$, $(5,1,26)$, $(5,-1.\,26)$ \\
$3,4;\ R3;\ N43$ & $(3,0,43)$, $(7,2,19)$, $(7,-2,19)$ \\
$3,4;\ N3;\ R43$ & $(6,3,23)$, $(11,5,14)$, $(11,-5,14)$
\end{tabular}
\end{center}
so zerfallen auch die negativen Klassen in vier Ordnungen:
\begin{center}
\begin{tabular}{l|l}
$3,4;\ N3;\ N43$ & $(-1,0,-129)$, $(-10,1,-13)$, $(-10,-1,-13)$ \\
$3,4;\ R3;\ R43$ & $(-2,1,-65)$, $(-5,1,-26)$, $(-5,-1,-26)$ \\
$1,4;\ N3;\ R43$ & $(-3,0,-43)$, $(-7,2,-19)$, $(-7,-2,-19)$ \\
$1,4;\ R3;\ N43$ & $(-6,3,-23)$, $(-11,5,-14)$, $(-11,-5,-14)$
\end{tabular}
\end{center}
Da aber das System der negativen Klassen dem der positiven stets so ähnlich ist, wird es meist überflüssig erscheinen, es gesondert aufzustellen. Wie sich die uneigentlich primitive Ordnung auf die eigentlich primitive zurückführen lässt, werden wir unten zeigen.

Für die Unterteilung der abgeleiteten Ordnungen brauchen wir schließlich keine neuen Regeln. Denn jede abgeleitete Ordnung stammt aus einer primitiven Ordnung (mit kleinerer Determinante), und ihre einzelnen Klassen entsprechen von selbst den einzelnen Klassen der primitiven Ordnung. Daher lässt sich die Unterteilung der abgeleiteten Ordnung offenbar aus der Unterteilung der primitiven Ordnung gewinnen.
